\documentclass[
uplatex,
b5paper,
10pt,
%english,
dvipdfmx
]{jsbook}
\usepackage{nruby}
\usepackage{type1cm} % 任意サイズの拡大縮小を可能にする
\usepackage[dvipdfmx]{hyperref}
\usepackage{pxjahyper}
\hypersetup{% hyperrefオプションリスト
 setpagesize=false,
 bookmarksnumbered=true,%
 bookmarksopen=true,%
 colorlinks=true,%
 linkcolor=black,
 urlcolor=black,
 citecolor=red,
}

\usepackage{framed}
\usepackage{wrapfig}
\usepackage{scalefnt}
\usepackage{version,url,here}	% required for `\comment' (yatex added)
\usepackage[dvipdfmx]{graphicx}	% required for `\includegraphics' (yatex added)
\usepackage{natbib,url}
\usepackage{pdfpages}
%\usepackage[varg]{txfonts}
\usepackage{makeidx}
%\usepackage{listings,jlisting}
\usepackage{listings,plistings}
\usepackage{color}
%\usepackage{xcolor}
\lstloadlanguages{[LaTeX]TeX, sh}
\colorlet{lstcolTeX}{green!50!black}
\colorlet{lstcoltext}{black}
\colorlet{lstcolshell}{blue!50!black}

\newcounter{marginparcntbw}[chapter]
\newcommand{\theMarginparcntbw}{$\dagger$\arabic{marginparcntbw}}
\newcommand{\Marginparbw}[2][−10pt]{%
  \stepcounter{marginparcntbw}%
  \textsuperscript{\theMarginparcntbw}%
  \protect\marginpar{\vskip#1\footnotesize%
    \textsuperscript{\theMarginparcntbw}
    {#2}\par}}


\setlength{\fboxsep}{.5zw}
\setlength{\fboxrule}{.6pt}

\begin{comment}
\lstset{basicstyle=\small\ttfamily, keywordstyle={}, commentstyle={},
  columns=flexible, showspaces=false, showstringspaces=false,
  aboveskip=12pt, belowskip=12pt, frame=tb,
  framesep=8pt, framerule=2pt, xleftmargin=10pt,
  xrightmargin=10pt, framexleftmargin=10pt, framexrightmargin=10pt
}
\end{comment}

\lstdefinestyle{shell}{language=sh, rulecolor=\color{lstcolshell!25}}
\lstdefinestyle{TeX}{language=TeX, rulecolor=\color{lstcolTeX!25}}
\lstdefinestyle{text}{language=TeX, rulecolor=\color{lstcoltext!25}}
\begin{comment}
\lstnewenvironment{listing}[1][]
{\lstset{#1}}
{}
\end{comment}


\lstset{% 
language={C++}, 
% backgroundcolor={\color[gray]{.85}},% 
basicstyle={\small},% 
identifierstyle={\small},% 
%commentstyle={\small\ttfamily \color[rgb]{0,0.5,0}},% 
%keywordstyle={\small\bfseries \color[rgb]{0,0,1}},% 
ndkeywordstyle={\small},% 
stringstyle={\small\ttfamily}, 
frame={tb}, 
breaklines=true, 
columns=[l]{fullflexible},% 
numbers=left,% 
xrightmargin=0zw,% 
xleftmargin=3zw,% 
numberstyle={\scriptsize},% 
stepnumber=1, 
numbersep=1zw,% 
morecomment=[l]{//}% 
} 

\begin{comment}
\lstset{
    basicstyle={\ttfamily\footnotesize}, %書体の指定
    keywordstyle={\footnotesize\ttfamily},%
%    frame=tRBl, %フレームの指定
    frame={t}, %フレームの指定
    framerule=1.2pt, %フレームの指定     
%    frame=leftbar, %フレームの指定
    framesep=0pt, %フレームと中身（コード）の間隔
    breaklines=true, %行が長くなった場合の改行
    linewidth=1\textwidth, %フレームの横幅
    xrightmargin=3zw,%
    xleftmargin=3zw,%
    numbers=left,%
    numberstyle={\ttfamily\scriptsize},%
    lineskip=0\baselineskip, %行間の調整
    tabsize=2 %Tabを何文字幅にするかの指定
}
\end{comment}

%\usepackage[format=hang,labelsep=colon,margin=10pt,sc,small]{caption}
\bibpunct[:\,]{(}{)}{,}{a}{}{,}
\definecolor{royalblue}{rgb}{0.0, 0.14, 0.4}
\newcommand{\colorrule}[1]{%
\begingroup\color{#1}\hrule\endgroup%
}%
\newcommand{\Colorrule}[1]{%
\begingroup\color{#1}\rule{1\textwidth}{2.4pt}\endgroup%
}%
\newcommand{\TColorrule}[1]{%
\begingroup\color{#1}\rule[2pt]{1\textwidth}{.6pt}\endgroup%
}%

\usepackage{framed}
\usepackage{color}
\usepackage{dcolumn}
\newcolumntype{d}[1]{D{.}{.}{#1}}
\usepackage{multicol}

\definecolor{lightgray}{rgb}{0.75,0.75,0.75}

\newtheorem{theo}{定理}[section]
\newtheorem{defi}{定義}[section]
\newtheorem{lemm}{補題}[section]

\makeatletter
\renewenvironment{leftbar}{%
%  \def\FrameCommand{\vrule width 3pt \hspace{10pt}}%  デフォルトの線の太さは3pt
  \def\FrameCommand{\vrule width 1pt \hspace{10pt}}% 
  \MakeFramed {\advance\hsize-\width \FrameRestore}}%
 {\endMakeFramed}
\makeatother

\newenvironment{redleftbar}{%
  \def\FrameCommand{\textcolor{red}{\vrule width 1pt} \hspace{10pt}}% 
  \MakeFramed {\advance\hsize-\width \FrameRestore}}%
 {\endMakeFramed}

\newenvironment{lightgrayleftbar}{%
  \def\FrameCommand{\textcolor{lightgray}{\vrule width .5zw} \hspace{10pt}}% 
  \MakeFramed {\advance\hsize-\width \FrameRestore}}%
{\endMakeFramed}


\AtBeginDvi{\special{papersize=\the\paperwidth,\the\paperheight}}

\newcommand{\mini}[2]{%
\setbox0=\hbox{\tt#1}\dp0=4pt%
\setbox1=\hbox{\tiny#2}\ht1=4pt\dp1=7pt%
\leavevmode\vtop{\offinterlineskip\box0\box1}}

\newif\ifVOLONE
\newif\ifVOLTWO
\newif\ifVOLTHREE

\VOLONEfalse
\VOLTWOfalse
\VOLTHREEtrue


% Japanese用の条件マクロ
\newif\ifJapanese
\newif\ifEnglish
%\Japanesefalse
\Japanesetrue   

% English用の条件マクロ
\ifJapanese
\renewcommand{\lstlistlistingname}{プログラム一覧}
\Englishfalse
\else
\Englishtrue
\fi

\newif\ifBLANK
%\BLANKtrue
\BLANKfalse

\newif\ifBIB
%\BIBtrue
\BIBfalse

\newif\ifINDEX
\INDEXtrue
%\INDEXfalse

\newif\ifOUT
%\OUTtrue
\OUTfalse

\newif\ifSHADOW
%\SHADOWtrue
\SHADOWfalse

\newif\ifSchedule
%\Scheduletrue
\Schedulefalse

\newif\ifBook
%\Booktrue
\Bookfalse

\newif\ifPreface
\Prefacetrue
%\Prefacefalse


\newif\ifCHSummary
%\CHSummarytrue
\CHSummaryfalse

\newif\ifPOSTSCRIPT
%\POSTSCRIPTtrue
\POSTSCRIPTfalse

\newif\ifWORKBOOK
\WORKBOOKtrue
%\WORKBOOKfalse


\newif\ifPLAYGAME
%\PLAYGAMEtrue
\PLAYGAMEfalse

\newif\ifVOCAB
%\VOCABtrue
\VOCABfalse

\newif\ifDEVELOPPER
%\DEVELOPPERtrue
\DEVELOPPERfalse

\ifWORKBOOK
\DEVELOPPERfalse
\Prefacetrue
\Booktrue
\CHSummaryfalse
\fi



\newcounter{excount}
\setcounter{excount}{0}
\newcounter{kdcount}
\setcounter{kdcount}{0}
\newcounter{ancount}
\setcounter{ancount}{0}
\newcounter{columncnt}
\setcounter{columncnt}{0}


\ifBook
\newenvironment{toiquestion}{%
\vspace{-1\baselineskip}
\noindent
\refstepcounter{excount}
\begin{quote}%
 \bfseries 
 \hspace{-2zw}問\theexcount\hspace{.7zw}%
}{%
\end{quote}
%\vspace{1.5\baselineskip}
%\vspace{-.5\baselineskip}
}
\else
\newenvironment{toiquestion}{%
\vspace{-1\baselineskip}
\noindent
\refstepcounter{excount}
\begin{quote}%
 \hspace{-2zw}問\theexcount\hspace{.7zw}%
}{%
\end{quote}
%\vspace{1\baselineskip}
%\vspace{-.5\baselineskip}
}
\fi

\newenvironment{toianswer}{%
%\vspace{-2\baselineskip}%
 \begin{quote}
  \parindent=1zw
  \hspace{.6zw}%
}{%
 \end{quote}
}

\setlength\unitlength{1pt}


\ifEnglish
\title{{\LARGE Linguistics}\\\protect\Colorrule{red}\\{\normalsize In very common daily life}}
\author{Hilofumi Yamamoto\\{\small\sc Ph.\,D. in Linguistics}}
\date{Tokyo Institute of Technology}
\else
\title{言語学\\
\vspace{-.5\baselineskip}
\Colorrule{red}\\\normalsize ことば研究史}
\author{山 元 啓 史\\{\small Ph.\,D.\,in Linguistics}}
\date{東京工業大学}
\fi

\makeindex
\begin{document}
\frontmatter
%\maketitle
\thispagestyle{empty}
\setlength\unitlength{1pt}
\begin{picture}(150,70)(70,565)  
% \put(172,182){\fbox{\includegraphics[trim=0 0 0 0,clip,width=0.465\hsize]{Sewing1898Edit.jpg}}}
% \put(162,182){\fbox{\includegraphics[trim=0 0 0 0,clip,width=0.52\hsize]{TheDifficultLesson1884.eps}}}
 \put( 0,475){\linethickness{0.4mm}\line(1,0){600}}
 \put( 0,173){\linethickness{0.4mm}\line(1,0){600}}
\ifEnglish
 \put(170,578){\scalefont{2.0}\bfseries Idiomatic Japanese}
 \put(145,555){\scalefont{1.5}\bfseries The Secret of Advanced Japanese}
 \put(230,530){\scalefont{1.8}\bfseries Volume 3}
 \put(210,500){\includegraphics[trim=10 45 10 45, clip, width=12mm]{./edx-logo.eps}}
 \put(245,502){\scalefont{1.5}\itshape\bfseries Tokyo\,Tech\,X}
 \put(100,484){\scalefont{1.5}\bfseries Let's learn Idiomatic Expressions of Japanese}
% \put(210,450){\scalefont{3.0} Workbook}
 \put(160,145){\scalefont{2.0}\bfseries Hilofumi Yamamoto}
 \put(215,130){\itshape\bfseries Ph.\,D.\,in Linguistics}
 \else
 \put(170,578){\scalefont{2.0}\bfseries Idiomatic Japanese}
 \put(145,555){\scalefont{1.5}\bfseries The Secret of Advanced Japanese}
 \put(230,530){\scalefont{1.8}\bfseries Volume 3}
% \put(210,500){\includegraphics[trim=10 45 10 45, clip, width=12mm]{./edx-logo.eps}}
% \put(245,502){\scalefont{1.5}\itshape\bfseries Tokyo\,Tech\,X}
 \put(100,484){\scalefont{1.5}\bfseries Let's learn Idiomatic Expressions of Japanese}
% \put(202,162){\scalefont{0.8}\bfseries やま}
% \put(239,162){\scalefont{0.8}\bfseries もと}
% \put(275,162){\scalefont{0.8}\bfseries ひろ}
% \put(312,162){\scalefont{0.8}\bfseries ふみ}
% \put(200,145){\scalefont{2.0}\bfseries 山　元　啓　史}

 \put(160,145){\scalefont{2.0}\bfseries Hilofumi Yamamoto} 
 \put(215,130){\itshape\bfseries Ph.\,D.\,in Linguistics}
\fi
\end{picture}
\newpage

\ifEnglish
Idiomatic Japanese, Volume 3
\else
日本語らしい日本語: 上級への道 Volume 3
\fi

\setlength\unitlength{1pt}
 \begin{picture}(0,130)(60,100)

%  \put(63,40){\fbox{\includegraphics[trim=0 0 0 0,clip,width=0.3\hsize]{TheDifficultLesson1884.eps}}}
%  \put(63,40){\fbox{\includegraphics[trim=0 0 0 0,clip,width=0.3\hsize]{Sewing1898Edit.jpg}}}  
%  \put(190,37){\includegraphics[width=5mm]{./Cc-public_domain_mark_white.eps}}
  
\ifEnglish
%  \put(55,20){\scalefont{1.0} Sewing}
%  \put(65,06){\scalefont{1.0} 1898 by William-Adolphe Bouguereau  (1825--1905)}
  
  %File:William-Adolphe Bouguereau (1825-1905) - Sewing (1898) Edit.jpg  
\else
%  \put(55,20){\scalefont{1.0} Sewing}
%  \put(65,06){\scalefont{1.0} 1898 by William-Adolphe Bouguereau  (1825–1905)}
  
\fi
 \end{picture}
% https://commons.wikimedia.org/wiki/File:Mary_Lemon_Waller_-_Spring_Voices_royalacademyillu1896roya_0070.jpg
\vfill

Hilofumi Yamamoto, Ph.\,D. in Linguistics, Tokyo Institute of Technology

\textcopyright\,Hilofumi Yamamoto, 2019


\ifPreface
\ifEnglish
\chapter*{Preface}
\else
\chapter*{はじめに}
\fi

\ifEnglish
This is one of the Japanese textbooks by edX MOOCs Tokyo Institute of Technology Online Education Development Room.
There is a knack to studying languages.
You should not stick to the details of it.
It is also necessary to get used to thinking that it does not matter.
No matter how much the structure of the bicycle you study, you will not be able to ride a bicycle.
There are many people who are riding a bicycle with no knowledge of the structure of the bicycle.
Just watching a bicycle will not allow you to ride a bicycle.
Anyway, if you do not ride a bicycle, you will not be able to ride.
Let's try it.

\else

本書はedX MOOCs 東京工業大学オンライン教育開発室による日本語の教科書である。
言語の勉強にはコツがあります。
詳細なところにこだわらないことである。
どうでも良いと考えることに慣れることも必要である。
自転車の構造をいくら勉強しても、自転車には乗れない。
自転車の構造を知らないで自転車に乗っている人はたくさんいる。
自転車を見ているだけでは、自転車に乗れない。
とにかく自転車に乗らなければ、乗れるようにはなれない。
やってみよう。
\fi

\vspace*{1\baselineskip}
\begin{flushright}
 \includegraphics[width=.2\hsize]{bicycle201801.pdf}

 \ifEnglish
 Hilofumi Yamamoto, Ph.\,D.\\
 Professor of Linguistics\\
 Tokyo Institute of Technology\\
 \else
 {\large 山 元 啓 史}\hspace*{3zw}
 
 {\small 東京工業大学教授\hspace*{2zw}}
\fi
\end{flushright}

%\end{document} % for checking the first page.


\mainmatter

\ifEnglish
\chapter{Getting Started}
\else
\chapter{さあ、はじめよう}
\fi

\includepdf[pages=32-36, scale=.84, trim=70 80 70 80, clip, offset=0mm -7mm, pagecommand={\thispagestyle{plain}}]{./idiom_compact.pdf}
\label{113010_11Dec18} % U2S1

%%%%%%%%%%%%%%%%%%%%%%%%%%%%%%%%%%%%%%%%%%%%%%
%%   Presession  %%%%%%%%%%%%%%%%%%%%%%%%%%%%%
%%%%%%%%%%%%%%%%%%%%%%%%%%%%%%%%%%%%%%%%%%%%%%


%%%%%%%%%%%%%%%%%%%%%%%%%%%%%%%%%%%%%%%%%%%%%%
%%   Dictation   %%%%%%%%%%%%%%%%%%%%%%%%%%%%%
%%%%%%%%%%%%%%%%%%%%%%%%%%%%%%%%%%%%%%%%%%%%%%

\ifEnglish
\chapter{Commentary}

\else
\chapter{解説編}

\fi

例文、解説、小練習、関連事項を掲載した解説編である。
練習編と同様、10の表現が１回分として使えるようになっている。

%\section*{日本語のイディオムと文型： 201 -- 210}
\begin{enumerate}
\setcounter{enumi}{200}

%201
\index{ものともせず@ものともせず}
\index{ものともしない@ものともしない}
\index{もの@もの}
\index{ず@ず}
\item 寒さを\underline{   }ともせず、彼は薄着で出かけていった。 

\begin{itemize}
% \itemsep=-4pt
\item[□] もの
\item[◆] 「ものともせず」、「ものともしない」... ``問題としない''、``恐
	  れない''という意味。「もの」の表現は、いろいろなイディオムがあ
	  るので、辞書を見ておくとよい。

\item[◆] 「もの」は物事、事物、者（人）、などなんでも対象となる。「もの
	  ともせず」は「対象物としないで」ということから、「全然問題にし
	  ないで」という意味になる。
\end{itemize}

\begin{itemize}
% \itemsep=-4pt
 \item 世の批判を物ともせずに信念を貫き通した。
      (Undaunted by criticism, he held fast to his opinions.)
       研究社『新英和・和英中辞典』より
\item MF中田\footnote{中田英寿(なかだひでとし): 日本の若手、天才的MF。サッ
      カーセリエAローマのMFで優勝。元全日本代表。2006年ワールドカップド
      イツ大会を機に引退。}は、昨晩帰国し、時差などものともせず、序盤か
      ら２アシスト２ゴールの大活躍で、セリエAの貫禄\footnote{かんろく
      (=presence, dignity)} を見せつけた。
\item 人間が耐えられない厳しい寒さや暑さをものともせず、真空中でも生き
       延びる ``スーパー生物''がいる。クマムシ\footnote{日本国内はもとよ
       り、北極に近いグリーンランドから南極大陸、ヒマラヤ山中や深さ百五
       十メートルの海底にまで生息している。}がそれである。
\end{itemize}

%202
\index{たんなる@単なる}
\index{すぎない@すぎない}
\item その選手にとって、１００勝という数字は\underline{   }通過点にすぎないだろう。 

\begin{itemize}
% \itemsep=-4pt
\item[□] 単なる
\item[◆] 単なる〜にすぎない、単なる〜のようなもの、単なる〜だけのもの
	  「単なる」特別のことではない。「ただの、ふつうの、一般的な」と
	  言い換えられる。名詞にかかる。「単なる(N)にすぎない」「単に(V)
	  したにすぎない」
\end{itemize}

\begin{itemize}
% \itemsep=-4pt
\item 彼は単なる「あほ」です。彼のいったことを気にする方がおかしいです。
\item それは単なる状況証拠にすぎない
\item 若者のおしゃれ感覚は抜群で、単に『まね』をするだけでなく、ファッショ
      ン的にも優れたものを見ることができる。
\item さて、サイトウという名字、この仕事を始めるまで私は「斎藤」も「斉藤」
      も同じで、単に斉は斎の略字だと思っていた。しかし、「斉」は「一斉に」
      と使うように、「ととのう」「ひとしい」という意味で、音はセイ。齊の
      略字。斎は「潔斎」と使うように、「つつしみ」「ものいみ」という意味
      で、音はサイ。齊と示を組み合わせた齋の略字。つまり斎と斉は別の字だっ
      た。『毎日新聞』より
\end{itemize}

\begin{itemize}
% \itemsep=-4pt
\item[☆] 比べましょう
\end{itemize}

\begin{itemize}
% \itemsep=-4pt
\item バスの乗降の新システムを単に試みているだけではない。
\item バスの乗降の新システムは単なる試みではない。(実験ではない)
\item その重要な問題をただ単に先送りしたにすぎない。
\item その重要な問題を単なる先送りにしたにすぎない。
\end{itemize}


%203
\index{など@など}
\index{とりたて@取り立て}
\index{とるにたらぬ@取るに足らぬ}
\index{とるにたりない@取るに足りない}
\item 若者言葉は貧弱だ、流行語など取るに\underline{   }と考えるのは老化した大人
      の偏見だ。

\begin{itemize}
% \itemsep=-4pt
\item[□] 足らない, 足らぬ
\item[◆] 取るに足らない[文末] / 取るに足らぬ N という形が多いが、「ぬ/
	  ない」どちらでも可能。「重要じゃない、どうでもいい、つまらない
	  こと」という意味。「足らない」の古語「足らぬ」は「取るに足らぬ
	  こと」のように現代語でも使われる。
\end{itemize}

\begin{itemize}
% \itemsep=-4pt
\item 誰がかいた絵かということは、取るに足りない問題と言ってよい。
\item ほとんど取るに足らない理由を口実に、連続リンチ殺人を繰り返していた。
\item 円に対してのドル安によるアメリカ国内の物価上昇などは、取るに足らな
      いものである。
\item 私自身のわずかな戦争体験を述べても、年上の人にはとるに足りないもの
      と受け取られ、若い世代には、はるかな過去の話でしかない。
\end{itemize}


%204
\index{かねをかける@金をかける}
\index{かける@かける}
\index{という@という}
\index{かねる@かねる}
\index{かねない@かねない}
\item 金をかければ、必ずよいものができるという意見があるが、私たちはこれ
      には賛成し\underline{   }。

\begin{itemize}
% \itemsep=-4pt
\item[□] かねる
\item[◆] V-masu かねる...''〜できない''、``しにくい''という意味。できな
	  い、という意味の間接的な表現。否定文ではないけれども否定の意味
	  があるので注意すること。否定型 V-masu かねない...否定形にする
	  と ``〜する危険がある''という意味になる。
\end{itemize}

\begin{itemize}
% \itemsep=-4pt
 \item[○] V-masu かねる
\end{itemize}
\begin{itemize}
% \itemsep=-4pt
 \item 遺憾ながら、ご依頼に応じかねます。
 \item お調べになりたい電話番号は、お名前だけでは、分かりかねます。
 \item レストランで子供の行儀があまりに悪いのを見るに見かねてしかってしまった。
\end{itemize}

\begin{itemize}
% \itemsep=-4pt
 \item[○] V-masu かねない
\end{itemize}
\begin{itemize}
% \itemsep=-4pt
\item 名産だからといって、たったひとつの果物ばかりを作り続けていたのでは、
      地元経済のアキレスけん\footnote{アキレス腱はギリシア神話の英雄アキ
      レスの名前に由来している。母親のテティスはアキレスを不死身にしよう
      と、不死身の水、ステュクス河に、彼をひたしたが、母は彼の踵をつかん
      でいたので、そこを水にひたすことを忘れてしまった。そのせいで、踵だ
      けが不死身にならなかった。アキレスはトロイ戦争に参加して活躍したが、
      矢が、アキレスの唯一の弱点であった踵に当って、死んでしまった。}に
      なりかねないということだ。
 \index{ばかり@ばかり}
\item 具体的な経済的基盤が確実でなければ、欧州通貨統合は絵に描いたモチに
      なりかねない
\item 子供がウェートレスにぶつかりかねない。
\end{itemize}

\begin{itemize}
% \itemsep=-4pt
\item[◆] 「〜をかねる」(兼ねる / 兼用 / 兼任 )
\end{itemize}
\begin{itemize}
% \itemsep=-4pt
 \item 小さい家なので、食事をするところと応接間を兼ねている。
 \item 急死のため、当分の間、首相が外相を兼ねることとなった。
\end{itemize}

%205
\index{という@という}
\index{でんぶんのひょうげん@伝聞の表現}
\index{なにひとつとして@何一つとして}
\index{だれひとりとして@誰一人として}
\item その彫刻家は、何\underline{   }として今までに完璧だと実感できる作品はなかっ
      たという。

\begin{itemize}
% \itemsep=-4pt
\item[□] 一つ
\item[◆] 何一つとして、誰一人として...ない→後ろに否定が来て、全然〜な
	  い、という意味になる。
\end{itemize}
\begin{itemize}
% \itemsep=-4pt
\item とにかく闇夜は何一つ物が見えない。
\item 容疑者のアリバイから判断して、犯罪らしい具体的な行動は何一つ浮かん
      でこない。
\item 遺伝子分析から、これらの結果を支持するような証拠は何一つ得られなかった。
\item 試合には勝ったものの、リーグ落ちが決定した\footnote{1999年、サッカー
      Jリーグで、浦和レッズは最終戦、延長戦で試合に勝ったが、その試合途
      中に他の試合で、ジェフ市原が先に90分で勝利を決めたため、得点差でJ1
      リーグからJ2リーグに降格することが決定した。}のでは、何一つとして
      勝利したと実感できず、涙を飲む結果となってしまった。
\item 「政治蚊共は巨悪を働いても誰一人ム所へ行かぬ」(原文のまま)『毎日新
      聞』より
\end{itemize}

%206
\index{みうごきがとれない@身動きが取れない}
\index{とれない@とれない}
\index{かのうのひょうげん@可能の表現}


\item その店は猫の額ほどの店で、客が五人も入れば身動きが\underline{   }なる。 

\begin{itemize}
% \itemsep=-4pt
\item[□]  取れなく / できなく
\item[◆] 「身動きが取れない」体が自由にならない、という意味だが、自分の
	  体に問題があるのではなく、部屋がせまい、人が多い、お金がない、
	  スケジュールがいっぱい、など周囲に原因がある場合に言う。
\end{itemize}

\begin{itemize}
% \itemsep=-4pt
\item 体を動かすこと→「満員で身動きもとれない(できない)車内」
\item 自由に行動すること→「借金で身動きとれない(できない)」
\end{itemize}

\begin{itemize}
% \itemsep=-4pt
\item[☆] 意味を考えましょう
\end{itemize}

\begin{itemize}
% \itemsep=-4pt
\item 議会や業界の圧力で身動きがとれない結果なのだ。
\item 山手線の上野—御徒町駅間は274％、身動きがまったく取れない状態だ。
\item 今や若者にとっては、携帯は、これなくして身動き一つとれない必需品なのだ。
\item 小渕首相は国民世論の渦のなかで身動きが取れない。
\item 通産省からの働き掛けが強力で、身動きが取れないことを認める。
\end{itemize}

\begin{itemize}
% \itemsep=-4pt
\item[◆] 「A が取れない」：A...時間、連絡、責任、場所、予約、対策
\end{itemize}
\begin{itemize}
% \itemsep=-4pt
\item 勉強の時間がとれないのではなく、勉強がやりたくないのです。
\item どうやってこの失敗の責任をとるのでしょう。
\item ここで花見をしたいという人が多く、春になると場所を取るのがたいへん
      だ。
\item 一度工場の中に入ったら、昼の休憩時間になるまで、連絡はとれませんと
      いわれた。
\end{itemize}

\index{ねこのひたいほど@猫の額ほど}
\index{ほど@ほど}
\begin{itemize}
% \itemsep=-4pt
\item[◆] 「猫の額ほど」..とてもせまいこと。なぜ「犬」ではなく「猫」なの
	  かわからない。「猫の額ほどの部屋」とはあまりいわない。
\end{itemize}

\begin{itemize}
% \itemsep=-4pt
\item 東京では、一生、汗して働いても猫の額ほどの土地も手に入らない。
\item うちの庭は猫の額ほどだが、誰にも売れない大切な理由がある。
\item その地主は、ゴミ消却炉の建設の噂が流れるとすぐ、猫の額ほどの広さで
      できるだけ多くの反対派の人に売り飛ばした。
\end{itemize}


%207
\index{やらねば@やらねば}
\index{やさき@矢先}
\index{たちあがる@立ち上がる}
\index{ふくごうどうし@複合動詞}
\index{ぎもんし@疑問詞}
\index{はんごのひょうげん@反語の表現}
\item 彼の死は「私が\underline{   }、誰がやる」といって、立ち上がった矢先の不幸
      であった。

\begin{itemize}
% \itemsep=-4pt
\item[□] やらねば
\item[◆] 〜ねば、誰がやる(反語：他には誰もいない)。「ねば」は「なければ」
	  と同じ。「ねば」は「ぬ」(203)の ``ば'' 形。もとは古語で、現代
	  語に残った形なので、引用句や改まった話、文中で使われることが多
	  い。
\end{itemize}

\begin{itemize}
% \itemsep=-4pt
\item 間にあわせねばと、必死でレポートを書いた。
\item 病気を早く治すためにも、食べたくなくても食べねば、とは思うが、舌の
      感覚が失われた人間には、気の遠くなるような努力であった。
\item モーニングコールを友人に頼んでおかねば、と考える私だった。
\item 「やらねば」という思いとは裏腹に身体が思うように動かない。
\end{itemize}

%208
\index{おもてもうらも@表も裏も}
\index{つくす@つくす}
\index{しりつくす@知り尽くす}
\index{つもり@つもり}
\index{つもりだった@つもりだった}
\index{まったく@まったく}
\index{きがつく@気がつく}
\item その選手は、プロ野球界の表も裏も知り\underline{   }いるつもりだったが、裏
      方で自分たちを支えている人々の存在にはまったく気がつかないでいた。

\begin{itemize}
% \itemsep=-4pt
\item[□] つくして
\item[◆] 〜つくす... ``すっかり〜てしまう''という意味。
\end{itemize}

\begin{itemize}
% \itemsep=-4pt
\item その研究は1960年代にすべてやり尽くされている。
\item 現代オリンピックの商業的問題点をすべて知りつくしている。
\item その店のメニューはすべて食べ尽くしてしまった。
\item 言うべきことはすべて言いつくした。
\item 彼は、銀行業界を知りつくしている。
\item 閉店のため値下げして商品を売りつくした。
\item 連休のために借りてきたビデオを全部見つくした。
\end{itemize}

%209
\index{いいことに@いいことに}
\index{こと@こと}
\index{あとをたたない@後をたたない}
\index{へいきで@平気で}
\item 黒い袋で中身が見えないのを\underline{   }ことに、危険物や非常識な汚物を平
      気で捨てる人が後をたたなかった。

\begin{itemize}
% \itemsep=-4pt
\item[□] いい
\item[◆] 〜をいいことに
\end{itemize}

\begin{itemize}
% \itemsep=-4pt
\item その学生は、先生が見ていないのをいいことに、毎回カンニングを繰り返
      していたそうだ。
\item 親がいないのをいいことに、子供はいつまでもテレビゲームに夢中だ。
\item 台風で電車が止まったのをいいことに、会社を休んでしまった。
\item 法で規制されていないのをいいことに40％もマージンをとって平気で商売
      をしている店がある。
\item 日本語が分からないのをいいことに、店の悪口を言いたいだけ言っていた
      ら、「そうですか。そんなにこの店、趣味がわるいですか？」と流暢な日
      本語で話しかけられてびっくりした経験がある。
\end{itemize}

%210
\index{として@として}
\index{よそ@よそ}
\index{をよそに@をよそに}
\index{ところ@ところ}
\item 世界遺産\footnote{世界遺産(せかいいさん)条約に基づく文化遺産として
      登録する制度。世界遺産委員会の本部はパリ。日本では、奈良・法隆寺、
      姫路城、西本願寺、広島原爆ドームが世界遺産条約のリストに登録されて
      いる。}としての保存を訴える被爆者の声を\underline{   }に、ヒロシマ
      \footnote{当時(1993年)、原爆ドームをユネスコの「世界遺産」に、との
      署名運動が広島で広がっていた。しかし、申請の権限を持つ文化庁は「原
      爆ドームは文化財ではない」と否定した。後世に伝える遺産と位置づける
      広島市と市民は「同じ戦争の遺産であるポーランド・アウシュビッツの強
      制収容所が登録されているのになぜ」、「国には原爆を投下した米国への
      配慮があるのではないか」との批判がマスコミに流れた。}を伝えるもの
      が消えてゆくところであった。
\begin{itemize}
% \itemsep=-4pt
\item[□] よそ
\item[◆] 〜をよそに...〜に全然関心を払わないで、〜に関係なく、という意味
\end{itemize}

\begin{itemize}
% \itemsep=-4pt
\item 外の雨天をよそに一回戦３試合の熱戦を行った。
\item 長引く不況をよそに、家庭用テレビゲームは、好調に売行きを伸ばした。
\item その国では、不況、凶作、災害などの国民の苦しみをよそに、民意を切り
      捨てる政治が長い間続いた。
\item 親の心配をよそに、娘たちは夜の渋谷で遊んでいる。
\item その俳優は、夫人の心配をよそに、「別宅持とうがアタシの勝手でしょ」
      と答えていた。
\item 湖畔では、氷点下の厳しい寒さをよそに露天ぶろがにぎわっている。
\item 両親の反対をよそに、一人エルサレムに来て宗教学校に通って、哲学を学
      んだ。
\end{itemize}
\end{enumerate}

%\section*{日本語のイディオムと文型： 211 -- 220}
\begin{enumerate}
\setcounter{enumi}{210}
\setcounter{footnote}{0}

%211
\index{ために@ために}
\index{みられる@みられる}
\index{みて@みて}
\index{かのうのひょうげん@可能の表現}

\item 中程度の豪雨であったが、短時間であったために各地で川が氾濫したと
      \underline{   }。

\begin{itemize}
% \itemsep=-4pt
\item[□] 見られる
\item[◆] 「〜と見られる」、「〜と見て」という表現は、``〜と考えられる''、
	  ``〜と考えて、予想して、''という意味で実際に見ているのではない。
	  助詞「と」を使う点に注意。
\item[◆] 類似の表現に「〜と見て、...」がある。意味はほぼ同じであるが、
	  文中で使われる。
\end{itemize}

\begin{itemize}
% \itemsep=-4pt
\item 就職活動にかかる費用が増えたのは、不況で会社訪問に交通費を支給する
      企業が減ったことなどが、主な原因と見られる。
\item その監督はチームをワールドシリーズ優勝に導くなどの活躍をしたが、自
      ら引退の理由は述べなかったが、一部では、健康上の問題が引退の主な理
      由と見られている。
\item アメリカン航空123便(ボーイング737)の機長から「米軍機と見られる戦闘
      機と異常接近（ニアミス）した」 との報告があった。
\item ウエストは年代が高くなるほど太っている。一般に63センチが標準と見ら
      れているが、30代では3割、40代では4人に1人で66センチの人数と同数、
      50、60代では69センチが最も多かった(4人に1人)。
\item 私の見るところでは = in my view.
\end{itemize}


\begin{itemize}
% \itemsep=-4pt
 \item[◆] 「氾濫」の意味は、川の水がいっぱいになって出てきて、被害が出
	    ること。この意味から、いっぱいに出てきて、迷惑でよくないもの
	    が出回ることの意味がある。
\end{itemize}

\begin{itemize}
% \itemsep=-4pt
 \item 利根川が氾濫した。
 \item 日本の本屋では、子どもに見せられないような雑誌が氾濫している。
 \item 日常の日本語の中には、外来語が氾濫している。
\end{itemize}


%212
\index{ては〜、〜ては〜@ては〜、〜ては〜}
\index{ていく@ていく}
\index{ては@ては}
\item  書い\underline{ A }消し、書い\underline{ B }消しして、作品を仕上げていった。

 \begin{itemize}
% \itemsep=-4pt
  \item[□] A ては; B ては
  \item[◆] 「A ては、B、A ては、B して、」...動作を繰り返す様子、対比し
	    て述べる表現。
 \end{itemize}

\begin{itemize}
% \itemsep=-4pt
 \item 喰っては、寝、喰っては、寝してさえいれば、すぐ太りますよ。
 \item 雪国の子どもでも、何回でも、滑べっては、転び、滑べっては、転びして、
       みんな滑べれるようになるんです。
\end{itemize}

\begin{itemize}
% \itemsep=-4pt
  \item[◆] 次の文は「ては」が一度しか出てこないが、繰り返しや対比の意味
	    がある。
\end{itemize}

\begin{itemize}
% \itemsep=-4pt
 \item その後もよく親の家を訪ねては小遣いなどをもらっていたという。
 \item 岩肌に爪（つめ）をたてては一歩ずつよじのぼるような確かさ。
 \item 高層ビル建築用超巨大ロボット(縦52メートル、横73メートル)。
       3日で1階分を組み立てては、油圧ジャッキで上昇を続け、まるでキノ
       コが成長しているよう。雨にぬれずに作業も出来る。
 \item イチロー\footnote{プロ野球選手。本名、鈴木一郎、オリックスブルー
       ウェイブの外野手。2000年より大リーグ、シアトルマリナーズに移籍し
       た。2007年のオールスターゲームでは日本人初のMVPに選ばれた。イチロー
       は日本だけでなくアメリカでも何度か首位打者のタイトルを手にしてい
       る。} は、守っては、ベストナインに選ばれるし、打っては、7年連続の
       首位打者のタイトルをとっており、現代野球の天才といわれている。
\end{itemize}

\begin{itemize}
% \itemsep=-4pt
  \item[☆] 比べてみましょう。
\end{itemize}

\begin{itemize}
% \itemsep=-4pt
 \item こう気が散っていては、書けない。
 \item こう暑くては、仕事にならない。
\end{itemize}

%213
\index{ても、〜ても@ても、〜ても}
\index{よう@よう}
\index{ようとおもう@ようと思う}
\index{のみほす@飲み干す}
\index{ふくごうどうし@複合動詞}
\index{ぬ@ぬ}
 \item 焼け酒で一升ビン\footnote{現在では、紙パックやペットボトルが使用
       されているが、昔は、酒や油の約1.8リットル入る大きなビンや缶が使わ
       れていた。}の酒を飲み干そうと思っても、日頃飲まぬ私では、飲ん
       \underline{   }飲ん\underline{   }なくならない。

 \begin{itemize}
% \itemsep=-4pt
  \item[□] でも
  \item[◆] 〜でも(ても)、〜でも(ても)、｛おわらない｜なくならない｝...
	    繰り返しても終わらない状態を述べる表現。
 \end{itemize}

\begin{itemize}
% \itemsep=-4pt
 \item ベルトをするのをわすれてしまっために、上げても上げてもスボンがず
       り落ちてきて、ステージではみんなの笑いものになってしまった。
 \item 行っても行っても終わりのない道。
 \item テーブルをふいてもふいても、汚されるので、いやになった。
 \item たくさんあるので、使っても使ってもなくならない。
 \item むしってもむしっても生えてくる草をむしることは、むなしい努力。
 \item 打たれても打たれても、なお耐え続けて勝利を手にしたボクサー。
\end{itemize}

%214
\index{ば@ば}
\index{ばいいのに@ばいいのに}
\index{たらいいのに@たらいいのに}
\index{たら@たら}
\index{うちあける@うちあける}
\index{ふくごうどうし@複合動詞}
 \item 早く打ち明ければ\underline{   }のに、まだ言えないでいる。

 \begin{itemize}
% \itemsep=-4pt
  \item[□]  いい
  \item[◆]  〜ば、いいのに／〜たら、いいのに
 \end{itemize}

\begin{itemize}
% \itemsep=-4pt
 \item 映画の内容を批判するなら堂々とやればいいのに。
 \item ちゃんと講義のノートを覚えればいいのに。
 \item 後ろから、ずっと白バイがついてきていて、気になって仕方がない。私
       は違反を犯していない自信はあるが、早く抜けばいいのについてくる。
 \item 「そんなに合格したいのなら、死ぬ気になってやればいいのに」とは言っ
       たが、本当に死んでしまうとは思わなかった。
 \item 「鬼のように」という表現も好きではない。「鬼のように勉強した」
       「鬼のように風邪を引いた」などと使われるので「すごく、とても」を
       表すらしい。それならはじめから「とても」や「すごく」を使えばいいのに。
\end{itemize}

%215
\index{ふくごうどうし@複合動詞}
\index{きる@きる}
\index{こむ@こむ}
\index{ところ@ところ}
\index{ところを@ところを}
\index{たところ、@たところ、}
 \item  疲れきって、ぐっすり寝込んでいる\underline{   }を起こされ、機嫌が悪くなった。

 \begin{itemize}
% \itemsep=-4pt
  \item[□] ところ
  \item[◆] 〜ているところを...、〜たところを...。「〜ところ\underline
	    {を}」のように助詞が「を」であることに注意。「ところ」は基本
	    的に「場所」という意味だが、「場面」という意味にもなって、
	    「ちょうど〜する/した/している/していた時」という意味になる。
 \end{itemize}

\begin{itemize}
% \itemsep=-4pt
 \item お忙しいところをお邪魔しました。
 \item 婚約者以外の女性といっしょにいるところを誰かに見られはしないかと
       心配していた。
 \item この生徒は、タバコを吸っているところを学校の担任の教師に見つかっ
       て、停学処分となっていた。
 \item １ページですむところを４ページも５ページも書くから、難しい文章に
       なるのだ。
 \item その犯人は自宅に戻ったところを、張り込んでいた刑事に逮捕された。
\end{itemize}

 \begin{itemize}
% \itemsep=-4pt
  \item[◆] 〜たところ、...「〜たら」の意味。
 \end{itemize}

\begin{itemize}
% \itemsep=-4pt
 \item 困っていたところ、助けてくれた。
 \item 間に合わないと思って急いで行ったところ、まだ電車はきていませんで
       した。
 \item カレーを作ったところ、みんなから、おいしいと喜んでもらえた。
\end{itemize}


%216
\index{のなにはじない@の名に恥じない}
\index{まぎれもなく@まぎれもなく}
 \item この二人の作家は、\underline{   }もなく中堅の実力派の名に恥じない小説家であっ
       た。 

 \begin{itemize}
% \itemsep=-4pt
  \item[□] まぎれ
  \item[◆] まぎれもなく...「まちがいなく」「疑いなく」「絶対に」という
	    意味。「まぎれる」という動詞があり、その意味は、「まちがえる」
	    という意味なので、「まぎれもなく」というのは「まちがいなく／
	    絶対に」という意味になる。その他、「気がまぎれる」という言葉
	    もあるが、これは「いやなことを忘れることができる」という意味。
	    「まぎれ\underline{が}なく」のような言い方はない。
 \end{itemize}

 \begin{itemize}
% \itemsep=-4pt
  \item 冠婚葬祭はまぎれもなく政治の場である。
  \item キリスト教との出合いは、まぎれもなく田中さんの人生を変えたのだ。
  \item ソウルの街頭でよく売っている、三十円くらいで中にはアンがいっぱい
	入っているたべもの。韓国語の名称は、タイ焼きではなく「ブンオ・パ
	ン」。「フナの形をしたパン」という意味。形や色はまぎれもなく、日
	本のタイ焼きと同じだが、韓国人は、タイよりフナに形が似ていると主
	張する。
 \end{itemize}

 \begin{itemize}
% \itemsep=-4pt
  \item[◆] 〜の名に恥じない...「〜といわれるくらい立派な」
 \end{itemize}

\begin{itemize}
% \itemsep=-4pt
 \item 横綱の名に恥じないよう、充実した毎日を送れるよう考えていきたいと
       思っています。
 \item 幸運に恵まれて伝統あるタイトルを獲得できました。今後はその名に恥
       じぬよう練習を続けたいと思います。
 \item 昨年１２月出版された「就職バカ本（ぼん）」（青英舎、９８０円）は
       編集部が学生の失敗談ばかり集めて、投稿とそれを題材にした４コママ
       ンガで構成したもので、題名に恥じない（？）「バカ話」が、いっぱい
       詰まっている。
 \index{ばかり@ばかり}
\end{itemize}

%217
\index{おもいおもいに@思い思いに}
 \item 三人の子供たちは、思い\underline{   }にこの一年の楽しかった想い出を日記に綴っ
       た。 

 \begin{itemize}
% \itemsep=-4pt
  \item[□] 思い
  \item[◆] 思い思いに、...「それぞれが、個人の好きなように」という意味。
 \end{itemize}

 \begin{itemize}
% \itemsep=-4pt
  \item 冬休み明けの月曜日。校内は、授業を終え、休暇の興奮を残した学生た
	ちが思い思いにくつろいだり、かたまって話し込んでいた。
  \item 仮に「公園体験」と名付けよう。ヨーロッパに旅行して、彼の地の公園
	で市民たちが思い思いにくつろぎ、遊んでいる光景に接して「うーん、
	いいなあ」と思う。
  \item 親子連れやカップルがウキワで沖合に出たり、砂浜で日光浴をするなど
	思い思いに浜辺の夏を楽しんでいた。
 \end{itemize}


%218
\index{あいついで@相次いで}
\index{という@という}
 \item 今場所は初日から二横綱が\underline{   }、やぶれるという波乱続きの幕開け
       であった。 

 \begin{itemize}
% \itemsep=-4pt
  \item[□] 相次いで
  \item[◆] 相次いで、〜という...続き。よくないことが続くことを書く場合
	    が多い。「次々に」とほぼ同じ意味。
 \end{itemize}

\begin{itemize}
% \itemsep=-4pt
 \item 相次いで、殺人事件が起こった。
 \item 相次いで、両親がこの世を去った。
 \item 英政府は八日夜、シェットランド諸島南部沿岸一帯を禁漁区としたが
       \footnote{1993年英スコットランド北部のシェットランド諸島沿岸でリ
       ベリア船籍タンカー「ブレイア」(89,700t)が座礁し、原油が流出し海上
       汚染事件がおこった。}、同諸島産のサーモンなどの契約キャンセルが相
       次いでいるという。
 \item 昨晩、名神高速道路下り線で、乗用車１台とトラック３台が相次いで追
       突、最後尾のトラックは横転した。
\end{itemize}

\begin{itemize}
 \item[比較] 次々と／に
\end{itemize}
\vspace{5mm}


%219
\index{ものの@ものの}
\index{たところ@たところ}
\index{ところ@ところ}
 \item 実験を数回繰り返したところ、着火はする\underline{   }すぐに火は消えることが分
       かった。 

 \begin{itemize}
% \itemsep=-4pt
  \item[□] ものの
  \item[◆] 「〜はするものの」→``〜するけれども、''という意味。
 \end{itemize}

\begin{itemize}
% \itemsep=-4pt
 \item 飛行機に乗って、日本に来たものの、言葉はわからないし、友達もいな
       いし、 自分の夢とはほど遠いスタートだった。
 \item 50代ではじめたハーフマラソン\footnote{普通のマラソン(フルマラソ
       ン)の42.195kmに対して、その半分の21.0975kmを走る。}では、アキレス
       けんが痛みだすハプニングがあったものの、夫婦２人そろって完走を果
       たした。
 \item 消費税の値上げはしないとは約束したものの、いつまでこの言葉を信じ
       ていいものか。
 \item マイケル・ジョーダン\footnote{Michael Jordan:元アメリカ、シカゴ、
       ブルズのバスケットプレーヤー。}という名前こそでていないものの、
       NBA\footnote{全米プロバスケットボール協会}の大物となれば、自ずと
       ジョーダン氏と誰もが考えてしまう。
\end{itemize}

%220 
\index{さき@先}
\index{さきほど@さきほど}
\index{さっき@さっき}
 \item この先どのようにするかは、\underline{   }に決定した法案を元に
      検討するとのことであった。 

 \begin{itemize}
% \itemsep=-4pt
  \item[□] 先
  \item[◆] つい、さきほど
 \end{itemize}

\begin{itemize}
% \itemsep=-4pt
 \item 先にゆでておいた卵をつぶして、今度はマヨネーズといっしょにソース
       を作ります。
 \item 雨なので、遠足は先に延ばしましょう。
 \item アインシュタイン\footnote{Einstein, Albert; 1879--1955; ドイツ生
       まれの米国の物理学者; 相対性理論の創始者}は、この先、相対性理論
       \footnote{そうたいせいりろん; the theory of relativity}を発表し、
       世界的な科学者になるのだが、当時、そんなことを想像する人は誰もい
       なかった\footnote{アインシュタインやエジソンは子どもの時、学校嫌
       いだったことが伝えられている。若林実著『アインシュタインも学校嫌
       いだった』『エジソンも不登校児だった』(筑摩書房)。}。
 \item[注] さっき、町で見た人は、ここの市長さんなんだよ。
\end{itemize}


 \begin{itemize}
% \itemsep=-4pt
  \item[【練習】] 「先」の意味は、\underline{過去のこと}、\underline{未
		  来のこと}、\underline{はじめに}の意味、\underline {場
		  所}の意味のうちのどれかを考えなさい。
 \end{itemize}

\begin{enumerate}\itemsep=-4pt
 \item 先に行ってください。
 \item 先に乗った人から、後ろに座ってください。
 \item 先に手順について説明しましょう。
 \item この先にきれいな喫茶店がありますから、行きましょう。
 \item 先のことは、誰にも分からないよ。
 \item 先日、お会いしたはずですが、名前を覚えていないんですよ。
\end{enumerate}
\end{enumerate}

%\section*{日本語のイディオムと文型： 221 -- 230}
\begin{enumerate}
\setcounter{enumi}{220}
\setcounter{footnote}{0}

%221
\index{つゆしらず@露知らず}
\index{ず@ず}
\index{ても@ても}
\index{かのうのひょうげん@可能の表現}
 \item そんな大変なことになっているとは\underline{   }知らず、家でゆっ
       くりビールでも飲んでいました。

 \begin{itemize}
 % \itemsep=-4pt
  \item[□] 露
  \item[◆] 「…とは露しらず」の形で覚えるとよい。全く気がつかなかった、
       知らなかったという意味。
 \end{itemize}

       「露」は昔から和歌でよく詠まれる（使用される）言葉（和歌でよく詠
       まれる言葉を歌語または歌ことばと呼ぶ）である。

       夜中に気温が下がって、早朝、草木についている小さな水玉のことを露
       という。気温が上がってしまうと、消えてなくなってしまうことから、
       「命のはかなさ(short-lived/ephemeral life)」をしばしば意味する。

 \begin{itemize}
 % \itemsep=-4pt
  \item 
 \end{itemize}

%222
\index{いてもたっても@いてもたっても}
\index{ず@ず}
\index{ても@ても}
\index{かのうのひょうげん@可能の表現}
 \item \underline{   }立ってもいられず。 

 \begin{itemize}
 % \itemsep=-4pt
  \item[□] いても
  \item[◆] 安心できない、落ち着かない様子、気持。
 \end{itemize}

 \begin{itemize}
 % \itemsep=-4pt
  \item 中学生の時、広島カープの初優勝に感動、それ以来、野球少年となった。
	大学でも野球を続けたが、ごく普通の会社に就職した。しかし、ある時、
	いても立ってもいられなくなり、安定した生活を捨て、プロ野球の入団
	テストを受けた。
%  \item 勝っても負けても一生懸命。
%  \item やせても枯れても一部上場企業。
%  \item 第二に勝っても負けても表情に出さない禁欲的なドラマ。
%  \item 家庭にいても外にいても、僕はいつも同じなんです。
%  \item 季語があってもなくても、17文字音であればよいのだ。
%  \item 競走馬は、重すぎても軽すぎても出世は難しいようだ。
 \end{itemize}

%223
\index{だめもと@だめもと}
 \item とても難しいとは思いますが、あきらめずに、
       \underline{   }でやってみましょう。

 \begin{itemize}
 % \itemsep=-4pt
  \item[□] だめもと
  \item[◆] 漢字で「駄目元」とも書く。「駄目で元々」の略。「元から駄目な
	    のだから、失敗しても惜しくはないはず。少ない可能性でも成功す
	    れば得だから、チャレンジした方がよい」の意味。
 \end{itemize}

 \begin{itemize}
 % \itemsep=-4pt
  \item 裕福に育った現代っ子は、何事も打算的で、だめでもともと、だめもと
	で、チャレンジするという姿勢が欠けていると言われている。
 \end{itemize}

%224
\index{とね@とね}
\index{みないとね@みないとね}
 \item いいとは思うけれど、高い品物なので買ってもいいかどうか、
       かみさんに聞いてみない\underline{   }。

 \begin{itemize}
 % \itemsep=-4pt
  \item[□] とね
  \item[◆] 「聞いてみないとね」は「聞いてみないとわからない」「聞いてみ
	    ないといけない」のように（わからない／いけない）という言葉が
	    省略されている。\footnote{1999年に所ジョージが作曲した「明石
	    家さんまさんに聞いてみないとね」（明石家さんま＆所ジョージ）
	    はこの言葉を有名にした。}
 \end{itemize}

 \begin{itemize}
 % \itemsep=-4pt
  \item 被災地が今どういう状況なのかわからないので、実際に行って
	みないとね、わかりませんね。
 \end{itemize}


%225
\index{こどもだまし@子供だまし}
\index{やすやすと@やすやすと}
\index{てにのる@手に乗る}
 \item そんなにやすやすと子供\underline{   }の手には乗るもんか。


 \begin{itemize}
 % \itemsep=-4pt
  \item[□] だまし
  \item[◆] 「子供だまし」簡単な人をだます手段。
       「やすやすと」簡単に。
       「手に乗る」相手の思うように動かされること。
 \end{itemize}

 \begin{itemize}
 % \itemsep=-4pt
  \item そもそもロックンロールそのものが「子供だまし」だからだ。誰でも弾
	けるような曲を弾いて、下手な声でがなる。変な格好で踊る。俺は強い
	ぜーと言ってみる。子供だまし以外の何者でもない。
       (\url{http://d.hatena.ne.jp/Maktab/20070501}より引用)
 \end{itemize}


%226
\index{いきをのむ@息をのむ}
\index{いきをこらす@息を凝らす}
\index{おもわず@思わず}
\index{じっと@じっと}
\index{かたずをのむ@固唾をのむ}
 \item その歴史的瞬間をじっと息を\underline{   }テレビで見ていました。

 \begin{itemize}
 % \itemsep=-4pt
  \item[□] のんで
  \item[◆] 「思わず、ハッと、じっと」などの言葉につづけて、驚いて息をと
	    めること。「息を殺して」「息を凝らして」「かたずをのんで」と
	    もいう。
 \end{itemize}

 \begin{itemize}
 % \itemsep=-4pt
  \item この紅葉の美しさに思わず息を飲んでしまった。
 \end{itemize}

%227
\index{いられない@いられない}
\index{ずにはいられない@〜ずにはいられない}
\index{にはいられない@〜（ず）にはいられない}
\index{ず@ず}
\index{かのうのひょうげん@可能の表現}
\index{くさい@〜くさい}
 \item 水着のデザインが古くさかったので、失礼だったかもしれな
       いが、笑わずには\underline{   }。


 \begin{itemize}
 % \itemsep=-4pt
  \item[□] いられなかった
  \item[◆] 「〜ずにはいられない」は「自然と〜してしまう」「〜したくなっ
	    てしまって困る」「〜したくてがまんできない」という意味。
 \end{itemize}

 \begin{itemize}
 % \itemsep=-4pt
  \item その漫才師のあまりのばかばかしさに、吹き出さずにはいられなかった。
  \item 愛さずにはいられない。
 \end{itemize}


%228
\index{せいいっぱい@精一杯}
\index{いっぱいいっぱい@いっぱいっぱい}
\index{なんて@〜なんて}
\index{こと@こと}
 \item 自分のことだけで\underline{   }一杯なのに、これ以上、他の人の
       世話なんてできません。

 \begin{itemize}
 % \itemsep=-4pt
  \item[□] 精
  \item[◆] 「精一杯やる」は、「一所懸命にやる」と文脈的には同じ意味もあ
	    るが、「これ以上できないほどの」という意味もあって、否定的に
	    使われることもある。
 \end{itemize}

 \begin{itemize}
 % \itemsep=-4pt
  \item 今はこれが精一杯、これで気に入らなければ、他の店に行って。
  \item 精一杯やって、それでだめだったら諦めるしかない。
  \item これでももういっぱいいっぱい。これ以上もう何も車に積み込めませんよ。
 \end{itemize}

%229
\index{もの@もの}
\index{すぎる@すぎる}
\index{かんがえもの@考え物}
 \item 子供が喜ぶのはいいことだが、こんなにはしゃぎすぎるも考え
       \underline{   }だね。
 \begin{itemize}
 % \itemsep=-4pt
  \item[□] 物
  \item[◆] 「考え物」は、「表面はよさそうだけれども、実は問題点」という
       意味。
 \end{itemize}

 \begin{itemize}
 % \itemsep=-4pt
  \item 裏側のデザインだけが派手になるのは考え物だ。
  \item 安いことはいいことだが、故障が少なからずあるのは考物だね。
  \item 人を使うことばかり覚えて、自分から何もしようとしなくなるのは考え物だ。
 \end{itemize}

%230
\index{まちにまった@待ちに待った}
\index{わくわく@わくわく}
\index{どきどき@ドキドキ}
 \item 今晩\underline{   }に待った初の世界陸上日本開催にわくわくドキ
       ドキしている。
       
 \begin{itemize}
 % \itemsep=-4pt
  \item[□] 待ち
  \item[◆] 「待ちに待った」は「いよいよ」とも言い替えられる。「長い間待っ
	    て、もうすぐその瞬間を迎える」という意味。

 \end{itemize}

 \begin{itemize}
 % \itemsep=-4pt
  \item 遅咲きの白いコブシの花が咲いた。樹齢250年のこの木はなぜか三年に
       一度しか花が開かず、今年が待ちに待った満開の年だった。
  \item 待ちに待った手紙がが今日届きました。
  \item 今日は待ちに待った海外旅行だったが、テロ対策のため空港のセキュリ
	ティチェックでちょ{}っとうんざりした。
 \end{itemize}

% 新しいアイテムはここに追加。

\end{enumerate}

%\section*{日本語のイディオムと文型： 231 -- 240}
\begin{enumerate}
\setcounter{enumi}{230}
\setcounter{footnote}{0}

%231
\index{よりけり@よりけり}
\index{ばあい@場合}
 \item それがいいかどうかは、時と場合に\underline{   }。 

 \begin{itemize}
% \itemsep=-4pt
  \item[□] よりけり
  \item[◆] 「けり」は古語では過去や詠嘆を意味するが、この場合は強調の意味。

 \end{itemize}

 \begin{itemize}
% \itemsep=-4pt
  \item 冗談をいうのは、時と場合、それに相手によりけりですよ。
  \item TPOというのは、Time、Place、Occasionの頭文字から作られた和製英語
	で、場合によりけり、という意味であったが、そのうち、条件に応じた
	正しい服装をする（Occasional Dressing）という意味になった。
 \end{itemize}

%232
\index{め@目}
 \item 彼には絵を見る\underline{   }がありますね。

 \begin{itemize}
% \itemsep=-4pt
  \item[□] 目
  \item[◆] 「目がある」は「人、芸術、価値の選択・判断する力がある」とい
      う意味。「選択眼がある」ともいう。
 \end{itemize}

 \begin{itemize}
% \itemsep=-4pt
  \item 人を見る目がある。
  \item 料理人の第一歩は、まず材料を見る目を養うことから始まる。
 \end{itemize}

%233
\index{ほかならない@〜に他ならない}
\index{なにをかくそう@何を隠そう}
\index{なに@何}
\index{とうぜんのごとく@当然のごとく}
\index{ごとし@ごとし}
 \item 物理学者が明らかにしようとしているものは、何を隠そうわれわれが当
       然のごとく生きているこの自然界に\underline{   }ならない。

 \begin{itemize}
% \itemsep=-4pt
  \item[□] 他
  \item[◆] 「〜に他ならない」は「〜そのものだ」に言い替えられる。「何を
      隠そう」は「本当のことをいえば」の意味。「当然のごとく」は「あたり
      まえのように」「日常何の不思議もなく」の意味。
 \end{itemize}

 \begin{itemize}
% \itemsep=-4pt
  \item なんでこんなにめんどうなことが毎日できるのかというとそれは楽しい
	からに他ならない。
  \item この選手の退場はチームにとって事実上の敗北宣言に他ならない。
 \end{itemize}

%ここにいれる。

\end{enumerate}

%\section*{日本語のイディオムと文型： 241 -- 250}

%\begin{enumerate}
%\setcounter{enumi}{240}
%\setcounter{footnote}{0}

%222
%\index{いてもたっても@いてもたっても}
%\index{ず@ず}
%\index{ても@ても}
%\index{かのう@可能}
% \item \underline{   }立っても居られず。 
%
% \begin{itemize}
%% \itemsep=-4pt
%  \item[□] いても
%  \item[◆] 
% \end{itemize}
%
% \begin{itemize}
%% \itemsep=-4pt
%  \item 今が大事なトマトの収穫期。しかし、「いても立ってもいられなくて」と
%	祖国での息子の晴
%  \item 勝っても負けても一生懸命。
%  \item やせても枯れても一部上場企業。
%  \item 第二に勝っても負けても表情に出さない禁欲的なドラマ。
%  \item 家庭にいても外にいても、僕はいつも同じなんです。
%  \item 季語があってもなくても、17文字音であればよいのだ。
%  \item 競走馬は、重すぎても軽すぎても出世は難しいようだ。
% \end{itemize}
%\end{enumerate}

%\chapter{基礎データ編}

%\input{kanji.txt}


\begin{comment}
\begin{enumerate}
 \item Unit 2 Section 1 p.\,\pageref{113010_11Dec18} %01
\end{enumerate}
\includepdf[pages=-, scale=.84, trim=70 50 70 34, clip, offset=0mm -7mm, pagecommand={\thispagestyle{plain}}]{./DICTEE01-20181009.pdf}
\label{113010_11Dec18} % U2S1
\end{comment}

%%%%%%%%%%%%%%%%%%%%%%%%%%%%%%%%%%%%%%%%%%%%
%%   Summary   %%%%%%%%%%%%%%%%%%%%%%%%%%%%%
%%%%%%%%%%%%%%%%%%%%%%%%%%%%%%%%%%%%%%%%%%%%
\ifCHSummary
\ifEnglish
\chapter{Summary}
\else
\chapter{まとめ}
\fi
\fi

\appendix
%\def\thechapter{Appendix \Alph{chapter}}
%%%%%%%%%%%%%%%%%%%%%%%%%%%%%%%%%%%%%%%%%%%%
%%   Appendix   %%%%%%%%%%%%%%%%%%%%%%%%%%%%
%%%%%%%%%%%%%%%%%%%%%%%%%%%%%%%%%%%%%%%%%%%%
 
\chapter{推薦するもの}

\chapter{英訳}

\begin{itemize}\small
%\itemsep=-4pt
\ifVOLONE
\item[1.] It only rains when you don't have an umbrella with you, so you can never count on the weather forecast.
\item[2.] It's cold again today, though not as cold as yesterday.
\item[3.] If I said that I'll do it, I'll see it through to the end.
\item[4.] It's locked, so they've probably gone out somewhere.
\item[5.] The weather is cloudy and grey, threatening to rain at any moment.
\item[6.] I can't figure it out at all. (I have no idea what's what.)
\item[7.] I looked out of the window for some reason or another and noticed that the leaves have turned red.
\item[8.] We are into late autumn, and the early mornings and evenings are quite chilly. How are you?
\item[9.] I read it after a fashion, but I didn't really understand what it was about.
\item[10.] Although he really knew, he pretended not to and told me nothing.
\item[11.] No matter how rapid the current is, the fish swim upstream.
\item[12.] She worked every day, without missing even a single day.
\item[13.] After you've arrived there safely, please give me a call right away.
\item[14.] Even though they say it's already spring, it's as cold as if winter were still here.
\item[15.] I have a slight headache, but since I have an exam today I have no choice but to go to school.
\item[16.] If you see Tanaka-san tomorrow, please tell him/her to give me a call.
\item[17.] One more failure and we will probably be compelled to call off the study.
\item[18.] No matter how much you eat, you don't necessarily gain weight.
\item[19.] The higher you climb on a mountain, the thinner the air becomes.
\item[20.] If I had to make a choice, I'd say that I like fish better than meat.
\item[21.] I looked over all the grammars I could lay my hands on, but no such thing appears anywhere in them.
\item[22.] After I conducted one study after another, the experiment finally yielded results.
\item[23.] Sunlight, water, and carbon dioxide are necessary for plants to grow.
\item[24.] I still can't get out of that childhood habit.
\item[25.] If I have to buy [something] regardless, I'd like to buy something that's cheap and good.
\item[26.] With this kind of academic record, you couldn't possibly get into a good university.
\item[27.] It's only natural for human beings to make mistakes.
\item[28.] People say there are too few parks in Japan, but is that really true?
\item[29.] This time very few people came, but we hope to have as many as possible next time.
\item[30.] Much to our regret, we are obliged to conclude that he is the criminal.
\item[31.] It is not until we have time to rest that we can lead our lives as human beings.
\item[32.] It's pointless to take the trouble to come to Japan and then not study Japanese.
\item[33.] I start to feel like travelling somewhere, though I don't have any particular place in mind.
\item[34.] On all the trains on the Tohoku line, skiers are packed in like sardines.(All the trains on the Tohoku line are jam-packed with skiers.)
\item[35.] I don't care who laughs at me.
\item[36.] Far from feeling relieved, I'm so worried that I can't sleep a wink at night.
\item[37.] It's been raining continuously since the day before yesterday.
\item[38.] On serene autumn nights, thoughts of my late mother come to mind.
\item[39.] The door, which should have been closed, is/was open.
\item[40.] It is rare to see such an interesting film.
\item[41.] Men are expected not just to be strong; men can't be men without gentleness.(=Real men can't just be strong; real men need to be gentle.)
\item[42.] I am told that that animal's body temperature rises along with an increase in the air temperature.
\item[43.] In the beginning it was easy, but gradually it has gotten more difficult.
\item[44.] There are just four of us here, so let's play mahjong.(=There are four of us here, which is just right for a game of mahjong.)
\item[45.] When you go to the post office, please stop by the supermarket and do the shopping.
\item[46.] The more you look at this picture, the more you can relish it(=the more interesting/pleasing it is to look at).
\item[47.] Customers are requested to note that tomorrow is our regular shop-holiday, so we will be closed.
\item[48.] I want to contact him straightaway, but unfortunately I don't have his phone number.
\item[49.] He always tells funny stories and makes us laugh.
\item[50.] A lot of trees burned down in the forest fire.
\item[51.] As long as we had water, we could survive for a few days.
\item[52.] No matter how hard it is, you must keep on trying right to the end.
\item[53.] The accident doesn't seem serious.
\item[54.] I'd like to show my appreciation to you with this little gift.
\item[55.] Although I thought it was a little expensive, I took the plunge and decided to buy it.
\item[56.] His attitude changed as suddenly as if he'd done a complete about-face.
\item[57.] People are queueing up and waiting to buy tickets.
\item[58.] The talk the teacher gave us yesterday attracted my interest.
\item[59.] The book I thought I'd lost turned up under the desk.
\item[60.] That's a bargain I found in a shop which I just happened to be passing by.
\item[61.] I inadvertently left my bag on the train.
\item[62.] The sight of the terrible traffic accident sent a chill down my spine.
\item[63.] In Taiwan, people wash their faces with cold water even in the winter.
\item[64.] Their children are already out on their own, so that couple lives quietly and comfortably.
\item[65.] We'd be very pleased to have you with us.
\item[66.] I cut my hand with a knife and it's covered with blood.
\item[67.] Fresh bread is very delicious.
\item[68.] You can't take the day off work just because it's going to rain.
\item[69.] In December, music associated with Christmas (=Christmas music) plays in town.
\item[70.] Although I know I shouldn't drink, I start having a good time and then I drink in spite of myself.
\item[71.] It appears unlikely to rain tomorrow because the sunset is [so] beautiful.
\item[72.] ``Where does your back itch?  Here?'' ``Yes, yes, there. Oh, that feels good!''
\item[73.] The relationship between language and culture is too strong to break.
\item[74.] We won't get anywhere asking every person for his/her opinion, so let's take a vote.
\item[75.] ``The strong win out in the end, don't they.'' ``Yes, such is life.''
\item[76.] Even though I said I didn't want to, I was forced into drinking.
\item[77.] That guy always says things that succeed in annoying me.
\item[78.] If someone says Egypt, we immediately call to mind the pyramids.
\item[79.] It is entirely possible that the Japanese race crossed over from the south in this way.(The theory that the Japanese race crossed over from the south in this way is entirely plausible.)
\item[80.] When I called my friend's house, a total stranger answered the phone.
\item[81.] Although this question was easy, not a single person was able to answer it.
\item[82.] Without having to think, ``I wonder what that was,'' I knew right away that the child's balloon had popped.
\item[83.] It's terribly unfortunate that that person passed away at the age of 35, in the prime of his/her working years.
\item[84.] Suzuki-san took leave of this world before his/her time, at the age of 45, and went to his/her eternal rest.
\item[85.] If I could spare the time and/or the money, I'd like to go on a trip or something, but reality is not that kind.
\item[86.] During my four years at university, I abandoned myself to pachinko and mahjong, and never studied seriously.
\item[87.] That person has already experienced that kind of thing, so he/she won't make the same mistake twice.
\item[88.] I caused my parents a great deal of worry because of this incident.
\item[89.] Knowing him, he has no common sense, so he might do such a thing.
\item[90.] From the lotus blossom emanated a fragrance so pleasant that it was beyond description.
\item[91.] Although I'd been planning on watching the 9 o'clock news, I got absorbed in my book and forgot all about it.
\item[92.] Although he/she had been fooling around every day up till then, all of a sudden he/she started studying seriously, and it was as if he/she had turned into a completely different person.
\item[93.] Because my father is an alcoholic, he can't go even a short time without a drink.
\item[94.] These days it's rare to see a whole family sit down together for a meal.
\item[95.] According to statistics, no more than 8% of all families sit down together for meals.
\item[96.] Today is Sunday, but it's raining and I don't have any money, so I have no alternative but to stay at home and sleep.
\item[97.] He couldn't possibly be a foreigner.
\item[98.] As might be expected of a student in the vocal music program, his/her speaking voice is pleasant too.
\item[99.] Because it's your fault, I think that you should apologize.
\item[100.] I don't even have a 10-yen coin, let alone a thousand-yen bill, so I can't possibly pay the money back.
\fi
\ifVOLTWO
\item[101.] It would be impossible to go out in such a fierce snowstorm, even if I tried.
\item[102.] Even if I complained to the likes of a [company] president such as him/her, he/she wouldn't listen.
\item[103.] He/she sent us one letter, and we haven't heard a word from him/her since.
\item[104.] It's a waste to leave this room unused.
\item[105.] His death will be mourned all the more deeply because he was the only painter born during the Meiji era.
\item[106.] My carelessness caused a child to be seriously injured.  (My carelessness inflicted serious injury on a child.)
\item[107.] Even though he had caught a cold, he came in first in the [long-distance] race.
\item[108.] If you're going to visit him/her in the hospital, wait until he/she gets a little better.
\item[109.]  Since then, nuclear power has come to attract attention as a new [form of] energy.
\item[110.] I almost left my umbrella on the train, you know.
\item[111.] I had a terrible time making up my mind whether to have soba (buckwheat noodles) or udon (wheat noodles), and I finally ended up having curry and rice.
\item[112.] On a hot summer day, nothing could be better than drinking cold beer in a beer garden.
\item[113.] Just what effect did these scientific-cultural phenomena have on human beings?
\item[114.] Whether you go or not, it doesn't matter to me.  (Whether you go or not has nothing to do with me.)
\item[115.] The bus is ready and waiting to leave, except that Yamada-san isn't here yet.
\item[116.] He turned to me with an expression verging on anger and said, "What brought you here?"
\item[117.] Even if the sky is cloudy, it isn't [bad] enough for you to take an umbrella with you.
\item[118.] I didn't marry Kyoko because I wanted her inheritance or anything.
\item[119.] A ``Thank you very much.  This is a gift to show my appreciation...''\\B ``That's really not necessary.''(=``Such a gift isn't at all necessary.'')\\A ``Then, could you at least tell me your name, please?''
\item[120.] The lost boy burst into tears the instant he saw his mother, who had come to pick him up.
\item[121.] I ran as hard as I could so as not to miss the train, and I made it.
\item[122.] You're not a child [anymore], and you can't go on eating kiddie lunches forever.
\item[123.] Although I'm studying very hard, my grades aren't improving the way they should.
\item[124.] I have an entrance examination tomorrow, so I don't think I'll get much sleep tonight.
\item[125.] That person appears to have already finished his meal and is having a cigarette with his coffee.
\item[126.] The opinion of the minority should be valued, but in the real world it tends to get ignored.
\item[127.] Now, at that dealer's, you can buy a Benz for 360,000 yen.  Of course, though, it doesn't have an engine.
\item[128.] Not only human societies but even monkey societies have a hierarchy of their own.
\item[129.] It would be hard enough to finish this job in three days, so telling us to finish it in half a day is absolutely crazy.
\item[130.] That person has such a look on his face, but it's not by any means because he's angry.  He always looks as if he's in a bad temper.
\item[131.] Believe it or not, this boy can read a book with more than three hundred pages in three minutes and then explain its contents!
\item[132.] After I was so kind to him/her, he/she could at least have said thank you.
\item[133.] Wouldn't it be a bit difficult to polish off three steaks and three plates of rice in an hour?
\item[134.] He is an idiot who can't be helped.
\item[135.] A ``You saved my daughter's life(=My daughter owe you her life.). I wouldn't feel right if I didn't do something to thank you.''\\B ``Oh, no, please don't bother.''
\item[136.] In your everyday life, you should refrain from behaviour that might invite misunderstanding.
\item[137.] The brothers' unseemly dispute over the inheritance continues.
\item[138.] Even if what you're saying is true, violence is wrong.
\item[139.] I hear that his illness is already past the point of curability, but let's pretend we don't know.
\item[140.] It isn't good for parents to butt in on children's quarrels.
\item[141.] It seems that autumn typhoons are frequently accompanied by heavy rain.
\item[142.] If he wants to cry, just let him cry.  He'll probably cry himself out in a little while.
\item[143.] For now, I'll just pay you 3,000 yen.  Please wait until the end of the month for the rest.
\item[144.] Some shops are friendly to customers, while others treat them coldly.
\item[145.] You know, if you scold a crying child thoughtlessly and excessively, it's not effective.
\item[146.] We'd rather buy a house than have to pay such high rent.
\item[147.] With the [signal] ``Ready, go!'' all of the competitors started off at once.
\item[148.] It is also conceivable that he stole it, but I don't want to think so.
\item[149.] Sorry to trouble you (=This may be a rather abrupt question), but could you tell me whether there is a public lavatory near here?
\item[150.] Examinations will be given class by class.
\item[151.] A child suddenly dashed out into the road, and I almost ran over and killed him/her.
\item[152.] It's not worth doing, at such a high risk.
\item[153.] Stop [your] dawdling, and hurry up and go to school.
\item[154.] Who do you want to tell this wonderful news to, first of all?
\item[155.] I'd like to continue with this work, even if it results in a big loss.
\item[156.] Whether you go on a trip or start a business of your own, spend this money on whatever you like.
\item[157.] Since I've already done this much to prepare, now I have no choice but to carry it out.
\item[158.] I've been very kind to him, but on top of that he even has the nerve to ask me for money...I'm at a loss for words.
\item[159.] Even though he/she is always borrowing money from me, he/she says, "Mum, what would you like for a present?"--which of course he/she doesn't really mean. 
\item[160.] I had thought that person was respected by everyone, but now I know I misjudged him/her completely.  What an awful thing he/she did!
\item[161.] This product is on offer at a bargain price, made possible only by direct delivery from the producer.
\item[162.] This is an ideal reference book for exam preparation.
\item[163.] Since he is trying so desperately to tell his story, it can't all be a pack of lies.
\item[164.] You've worked so hard up till now, and now you're going to quit?  Don't be silly!
\item[165.] Even if it's "all-you-can-eat," you'll probably only be able to finish up three plates at the very most, no matter how much you eat.
\item[166.] If goods are priced too highly, they won't sell well.
\item[167.] There's such a big difference of opinion between us that it's not even worth mentioning.
\item[168.] No matter how hard I try, I'll end up being called names.
\item[169.] Because nobody was kind enough to do it for me, I had no choice but to do it myself.
\item[170.] It seems that the project isn't proceeding according to plan because of insufficient funds, and because of outside interference as well.
\item[171.] When all is said and done, in the end you can trust only yourself.
\item[172.] I'm not really sure, but I think the conversation was something about the question of pensions...
\item[173.] I can tell you, it was so hugely amusing that I almost died laughing.
\item[174.] I hear they were discussing so important a matter alone together, in whispers!
\item[175.] This politician's scheme met with the criticism of the opposition parties and, furthermore, provoked the antipathy of the entire nation.
\item[176.] Cancer is not the sort of illness which people catch all of a sudden; it's a complicated and incurable disease which gnaws away day by day at the health of the patient.
\item[177.] In winter, this place is a lonesome and deserted town, but in the summer, when the swimming season opens, it livens up with young people in colourful swimsuits.
\item[178.] At Hyde Park in England, strangers engage each other in lively discussions every day.
\item[179.] The whole mountain turned into a sea of fire, all because of the tiny flame of a cigarette.
\item[180.] A wooden house, for instance, is reduced to ashes in an instant once it catches fire.
\item[181.] Certainly I'd heard it was awful, but I really didn't expect it to be this horrendous.
\item[182.] Although I did it because I thought it was a good thing to do, it ended up making that person all the more unhappy.
\item[183.] For the next test, I'm just going to write down whatever any old book says and fake my way through.
\item[184.] There are more and more suicides among the children of parents who expect too much of them.
\item[185.] This cafe has many customers who come to savour not only the smell of its coffee but also the atmosphere of the place.
\item[186.] It was not until our baby was born that I actually started to feel like a father (=that the actual feeling of being a father welled up in me).
\item[187.] It's not that I don't cook at all, but when I'm alone it's a nuisance so I usually eat out.
\item[188.] I hear that person has gone to Tokyo and doesn't live here any more.
\item[189.] He came into my room silently, without even being invited.
\item[190.] Purely and simply because I studied very hard, I was able to get a good grade.
\item[191.] What with all the surprise and rejoicing, the place was in an uproar.
\item[192.] I can't get things sorted out unless I deal with them one at a time.
\item[193.] I can't figure out what's going on.  That person turned up with a year-end gift ("oseibo")!
\item[194.] When you grow up, you must not become the kind of person who causes nuisance to others, even by mistake.
\item[195.] That person is no prime minister, after he/she went so far as to embezzle public money.
\item[196.] Carried away by her sorrow, she couldn't even cry; she just squatted down without any expression on her face.
\item[197.] [The government] lowered income taxes, but meanwhile it raised the consumption tax, in order to guarantee its financial resources.
\item[198.] Now that things have come to this, we have no choice but to do a drastic restructuring and reorganize management.
\item[199.] I put a half-finished cup of coffee beside the computer while I was working, but I accidentally spilled it.
\item[200.] It's hard to believe that they can hunt down criminals by means of extrasensory perception.
\fi
\ifVOLTHREE
\item[201.] Making nothing of the cold, he went out in thin clothes.
\item[202.] The record of hundred wins would be merely a passing point for the player(pitcher).
\item[203.] It is a prejudice of the old to think that the language of the young is poor, and vogue-words are insignificant.
\item[204.] There is an opinion which we can always make a better product if we spend more money, but we don't agree with this.
\item[205.] The sculptor said that he had never had any perfect works which could satisfy him.
\item[206.] The shop is so tiny that it becomes impossible to move with only five people inside the shop.
\item[207.] His death was a tragedy happened just as he was about to take action after saying, ``who else would do it, if I don't.''
\item[208.] The player believed that he knew everything about the world of professional baseball, but he didn't realise the existence of the people supporting players behind the scenes.
\item[209.] Many people disposed garbage that contains dangerous and unacceptable goods taking advantage of black garbage bags that do not show their contents.
\item[210.] In spite of the voices appealing for the preservation of Hiroshima Peace Memorial as a world heritage, the landmark conveying the history of Hiroshima was about to be vanished.
\item[211.] It was a heavy rain but not a serious one. However we judge that it caused floods in the various places due to its short duration.
\item[212.] S/he finished a work up by drawing and erasing it again and again. After many repeats of writing and deleting, s/he completed the work.
\item[213.] I feel abandoned and try to drink up a 1.8-litter-bottle of sake, but I have had a lot of it yet even if I keep drinking since I don't usually drink.
\item[214.] Even though he should speak his mind, he is still hesitating about doing so.
\item[215.] He became angry since someone woke him up while he was sound asleep, exhausted.
\item[216.] These two writers are the novelists unmistakably appropriate to call the competent and influential. 
\item[217.] The three children wrote the pleasant memories of the year one by one in their diaries.
\item[218.] This sumo tournament has started with eventful episodes such as both of two Yokozunas being defeated one after the other on the first day.
\item[219.] After having conducted the experiments several times, it was found that although it could be ignited, the flame would  go out immediately.
\item[220.] What we should do in future will be discussed based on the bill which was legislated recently.
\item[221.] I was just drinking beer at home without realising that such a terrible thing happened.
\item[222.] I went onto a flight right away, feeling restless, when I heard of my father's accident.
\item[223.] I think that it is very difficult, but don't give up and go for it! You have nothing to lose.
\item[224.] I think it is quite good but because it is an expensive item, I'd better ask my wife whether it is okay for me to buy it.
\item[225.] I won't be easily cheated by such a childish trick.
\item[226.] I was watching the historic moment on TV with bated breath.
\item[227.] It may have been impolite, but I could not help laughing because the design of her swimwear was very old-fashioned.
\item[228.] I have too much on my plate already.I can't do anything else for others.
\item[229.] It is good that children are delighted, but they frolic too much that makes you wonder.
\item[230.] Tonight there will be a long waited game of IAAF(International Association of Athletics Federations), for the first time in Japan, and I am very excited about it.
\item[231.] It depends on circumstances whether it is proper or not.
\item[232.] He has a good eye for arts.
\item[233.] What physicists have been trying to unveil is, to tell the truth,  nothing else but something about this natural world where we live as a matter of course.
\item[234.] A politician who said last week ``I will make the politics of Japan better'' has been arrested today due to a corruption case.
\item[235.] ``About one person is so demanding that I did what she requested unwillingly.'' ``Do yo mean it's me?''
\item[236.] Looking at a wedding as an outsider, it seems a full of happiness, but for the couple themselves it is indeed a very troublesome event.
\item[237.] He is not at all such a weak person who would whinge in the presence of other people.
\item[238.] Saving the face of their Prime Minister cannot be the purpose of the politics of any country.
\item[239.] We feel a little awkward because we have just had our wedding ceremony but let us inform you that yesterday we found out that my wife is pregnant.
\item[240.] Yesterday, I ran into my dear friend at a most unexpected place.
\item[241.] My company went bankrupt, I got into a large debt, I got a divorce, I did't know what I to do, and I was at a loss.
\fi
\end{itemize}

%\chapter{プログラム編}

%\input{idiom.pl}



\backmatter
\ifPOSTSCRIPT

%%%%%%%%%%%%%%%%%%%%%%%%%%%%%%%%%%%%%%%%%%%%
%%   Postscript   %%%%%%%%%%%%%%%%%%%%%%%%%%
%%%%%%%%%%%%%%%%%%%%%%%%%%%%%%%%%%%%%%%%%%%%

\ifEnglish
\chapter*{Postscript}
\else
\chapter*{あとがき}
\fi

\section*{読むということ}

\vspace*{1\baselineskip}

\begin{flushright}
 {\large 山 元 啓 史}\hspace*{3zw}
 
 {\small 東京工業大学教授\hspace*{2zw}}
\end{flushright}

\begin{comment}
\begin{enumerate}
 \begin{comment}

 \item 授業は隣の人とのおしゃべりが中心、
 \item 毎回振り返りのレポート、
 \item 試験問題を学生自らが作り、
 \item 鉛筆を転がすような試験ではなく、
 \item 本質が議論できるよう口述試験を受け、
 \item テーマを自分で選び、
 \item 決められた作法で書くことを学び、
 \item 自分で何かを明らかにする、
 \item 徹底的に能動的でなければ参加できない形式にした。
 \end{comment}
\end{enumerate}

\end{comment}
\fi

\ifBIB
\bibliographystyle{jecon}
\bibliography{koten}
\fi

\ifINDEX
\newpage
\ifBook
\printindex
\else
\printindex
\fi
\fi

\cleardoublepage
\thispagestyle{empty}
\noindent
\ifEnglish
{\Large \underline{Back cover:}}
\else
{\Large \underline{裏表紙について}}
\fi
\setlength\unitlength{1mm}

\ifBook
 \begin{picture}(50,180)(0,0)
\else
 \begin{picture}(50,180)(0,0)
\fi   
% \put(80,117){\includegraphics[width=30mm]{../phd/p20/Quintus_Horatius_Flaccus.eps}}

\ifEnglish
% \put(80,110) {
% \begin{minipage}[c]{70mm}\scriptsize
%  Quintus Horatius Flaccus\\
%  BC.65.12.8--BC.8.11.27\\
%  Roman poet from Southern Italy
% \end{minipage}}
\else
%  \put(80,106) {
% \begin{minipage}[c]{70mm}\scriptsize
%  クィントゥス・ホラティウス・フラックス\\
%  Quintus Horatius Flaccus\\  
%  BC.65.12.8--BC.8.11.27\\
%  古代ローマ時代の南イタリアの詩人
% \end{minipage}}
\fi  
% \put(10,130){\fbox{\includegraphics[width=60mm]{../phd/p20/Carpe_Diem.eps}}} 

\ifEnglish
 \put(10,108){
% \begin{minipage}[c]{60mm}\small
%  In the sundial of the photograph, it is engraved with ``Carpe Diem'' in Latin, ``Seize the day'' in English.
%  It means that you will enjoy the day and living the best.
%  It is a phrase that appears in the poetry of ancient Roman poet Horatius in the 1st century BC.
%  The phrase also appears in the movie ``Dead Poets Society'', 1989 played by Robin William.
% \end{minipage}}
\else
% \put(10,100){
% \begin{minipage}[c]{60mm}\small
%   写真の日時計にはラテン語で``Carpe Diem''（カルペ・ディアム）と彫ってあります。
%  英語では``Seize the day''、
%  日本語では「その日を摘め」と訳されています。
%  そこには「その日を楽しみ、精一杯いきること」という意味があります。
%  紀元前1世紀の古代ローマの詩人ホラティウスの詩に登場する句で、
%  映画``Dead Poets Society''（1989年、邦題「いまを生きる」ロビン・ウィリアムズ主演）にも出てきます。
% \end{minipage}}
\fi  
  \put(20,0){
  \begin{minipage}[c]{100mm}\small
    \begin{center}
\ifBook
     Ten Sentences A Day for Eight Weeks\\Dictation Everyday\\Volume 3\\
\else
     Workbook\\    
\fi
\ifEnglish
     \textcopyright\ First edition\\
     
%     Author: Hilofumi Yamamoto\\
\else
     \today\  第1版\\
     
     著者: 山元啓史\\
\fi     
    \textcopyright 2018, Hilofumi Yamamoto
    \end{center}
  \end{minipage}}
  \end{picture}

\newpage
\thispagestyle{empty}
\ifBook
\begin{picture}(50,160)(0,0)\centering
\else
\begin{picture}(50,160)(0,0)\centering
\fi  
% \put(04,127){\includegraphics[bb=0 0 870 129,clip,width=106mm]{../hiratoki.sub_header.jpg}}
% \put(07, 51){\fbox{\includegraphics[width=100mm]{../phd/p20/Carpe_Diem.eps}}}
 \put(42, 41){\Large\sc Carpe Diem}
% \put(38, 30){\fcolorbox{white}{royalblue}{\includegraphics[width=35mm]{../phd/logo02.png}}}
% \put(50, 20){\includegraphics[trim=0 0 0 0, clip, width=10mm]{./edx.eps}}
% \put(100, 20){\includegraphics[trim=0 0 0 0, clip, width=10mm]{./oedo2.eps}}
% \put(36, 10){\includegraphics[trim=0 0 0 0, clip, width=40mm]{./maxresdefault.eps}}
% \put(33, 13){\includegraphics[trim=10 10 10 10, clip, width=23mm]{./edx-logo.eps}}
% \put(56, 10){\includegraphics[trim=10 100 10 150, clip, width=28mm]{./oedofblogo.eps}}
% \put(27, 15){\includegraphics[trim=185 250 185 250, clip, width=60mm]{../phd/hisaotakibanner2.pdf}}
% \put(37,-05){\includegraphics[width=40mm]{../phd/p20/guru_logo1.eps}} 
% \put(45.5, 3.5){\includegraphics[bb=0 0 421 178,clip,width=23mm]{../phd/KAKENHIlogo_M.jpg}}
% \put(38, -6){\includegraphics[trim=0 0 0 0, clip, width=38.5mm]{../phd/p20/KAKENHIlogo_M.eps}} 
\begin{comment}
 % \put(04,127){\includegraphics[bb=0 0 870 129,clip,width=106mm]{../hiratoki.sub_header.jpg}}
% \put(07, 51){\fbox{\includegraphics[width=100mm]{../phd/p20/Carpe_Diem.eps}}}
 \put(42, 41){\Large\sc Carpe Diem}
% \put(38, 30){\fcolorbox{white}{royalblue}{\includegraphics[width=35mm]{../phd/logo02.png}}}
% \put(42, 15){\includegraphics[bb=0 0 421 178,clip,width=30mm]{../phd/KAKENHIlogo_M.jpg}} 
% \put(37,-05){\includegraphics[width=40mm]{../phd/p20/guru_logo1.eps}}
\end{comment} 
\end{picture}
\end{document}

\begin{comment}
\section{ポスター発表とコメント（必修）}

\begin{enumerate}

 \item ポスター発表のタイトルの申し込みをGoogle Formから、7月20日の
       23:59:59（厳守: これを守らないと単位喪失）までにお願いします。
 \item ポスター発表のタイトルの申し込みは{\Huge こちら$\downarrow$}

       \begin{center}
       \includegraphics[trim=0 0 0 0, width=.3\textwidth, clip]{../phd/p20/LinguisticsB_PosterPresentation.eps}	
       \end{center}
\end{enumerate}

\end{comment}
